\documentclass[12pt, titlepage]{article}
\usepackage[margin=1in]{geometry}
\usepackage{graphicx}
\usepackage{float}

%\usepackage{hyperref}

\title{\Huge Wingfeather Flight Controls Lab Drone \\ \Large{User Manual}}
\author{by \\ \Large TEAM WINGFEATHER \\ Aydan Vissing \\ Brady Schick \\ Gabriel Southwell \\ Simon Ross}
\date{\footnotesize{Last Updated} \\ \normalsize{\today}}


%\setlength{\topmargin}{-.25in}
%\setlength{\oddsidemargin}{.25in}
%\setlength{\evensidemargin}{.25in}
%\setlength{\textwidth}{6.0in}
%\setlength{\textheight}{9.0in}
%\renewcommand{\baselinestretch}{1.4}
%\usepackage{indentfirst} %Maybe keep this maybe not

\begin{document}
\maketitle

%FIXME remove comments

%FIXME improve abstract

\newpage
\tableofcontents
%\listoffigures
%\listoftables

\newpage
\section{Quick Start Guide}

The user manual must include a Quick Start guide which should be a very simple one-page explanation of how to get the device set up for a first test using built-in default settings.

You must explain how to use all your instrument’s features. Assume that the reader of your user’s manual has never seen such an instrument. Your user manual will be their only reference source as they try to set up (i.e. connect the wires) and make use of the device. You may want to test your user manual on someone who has no experience like a friend from a different major. You may assume that the general user has a basic understanding of the terminology common to the project.

\section{Drone Construction}

\subsection{Ordering Components}
\subsubsection{PCB}

\subsubsection{Motors}
(key specs and a link)

\subsubsection{Propellers}
(key specs and a link)

\subsection{Printing the Frame}

\subsection{Assembling the components}

\subsection{Programming the Drone}

\subsection{Using the App}

\subsection{Using ESP IDF}

\section{Drone Development}

\subsection{Basic Repository Structure}

\subsection{Basic Drone Testing}

\subsection{PID Tuning}

\subsection{Frame Development}

\subsection{PCB Development}

\end{document}
